\begin{helper}
Assume
\(\noderef{ParameterHierarchyEventualComparisons}
(\varepsilon,\varepsilon_1,\varepsilon_2,n_0)\), \(n>n_0\), and
\[
k=\noderef{TwoBiteNaturalIndependenceScale}(1+\varepsilon,n).
\]
Let
\[
t_1=\noderef{TwoBiteHugeCutoff}(n),\qquad C=100(\log n)^3.
\]
Then \(t_1>0\), \(C\ge0\), and
\[
\binom{2k/t_1+1}{2}_{\mathbb R}C\le \frac12 k.
\]
\end{helper}
\begin{proof}
Apply \(\noderef{ParameterHierarchyEventualComparisons}\) at the present
\(n\).  One of its conjuncts is exactly
\[
\binom{2K/t_1+1}{2}_{\mathbb R}100(\log n)^3\le\frac12K,
\]
where \(K=\noderef{TwoBiteNaturalIndependenceScale}(1+\varepsilon,n)\).
The hypothesis \(k=K\) rewrites this as the displayed inequality with
\(C=100(\log n)^3\).

It remains to record the signs needed by the counting argument.  The same
eventual-comparisons package gives
\[
1\le 2pm,\qquad 0\le\varepsilon_2,\qquad
2pm\le(\varepsilon_2/100)t_1.
\]
Since \(\varepsilon_2/100\ge0\), the inequalities
\(1\le 2pm\le(\varepsilon_2/100)t_1\) force \(t_1>0\); otherwise the right
side would be nonpositive.  They also force
\[
p=\frac12\sqrt{\frac{\log n}{n}}>0,
\]
because \(m>0\) is another conjunct of the same package.  Since \(n>0\), the
positivity of this square root gives \(\log n>0\).  Therefore
\((\log n)^3\ge0\), and multiplying by \(100>0\) gives \(C\ge0\).
\end{proof}
